\section{Certified cyclic Garside preprocessing}
\label{sec:cyclic-garside}

Report 25 extends the earlier rank-two interval optimizer to arbitrary
supports in an Artin braid group.  The maintained implementation shares
exact prefix arithmetic across cyclic cuts and target words.  It supplies
an optional source-braid probe, not a new knot invariant.  A successful
proposal must shorten the current simplified diagram, must pass independent
replay, and restarts recognition under the remaining global budget.
The probe follows the existing braid, Seifert, descending, and modular
Alexander obstructions; its default setting is disabled.

\subsection{Exact states and a separately checked equality witness}

Let $b$ be the strand count and $\Delta$ the positive half-twist of length
$B=\binom b2$.  Classical Garside theory gives a unique left-greedy form
\[
 g=\Delta^z s_1\cdots s_r,
\]
where $z\in\mathbb Z$ and the $s_i$ are proper nonidentity simple braids.
A simple braid crosses each pair at most once and is determined by its
permutation.  These are classical facts, not new compression theorems;
see Dehornoy~\cite{dehornoy-garside}.  The implementation stores complete
permutation arrays and the integer power.  Sorting and binary search
compare these exact states; neither a permutation quotient nor a hash
fingerprint is used to decide braid equality.

Our convention is $p[i]=$ bottom position of top strand $i$.  For two
adjacent simple factors $A,C$, a positive atom $\sigma_{i+1}$ can move
from the start of $C$ to the end of $A$ exactly when
\[
 C[i]>C[i+1],\qquad A^{-1}[i]<A^{-1}[i+1].
\]
The move swaps values $i,i+1$ in $A$ and positions $i,i+1$ in $C$.
The first condition factors $C=\sigma_{i+1}C'$ and reduces simple length
by one; the second makes $A\sigma_{i+1}$ simple and increases its length
by one.  Thus the local product is unchanged.  An identity right factor
is deleted.  Positive input letters append an atom; negative letters use
\[
 \sigma_i^{-1}=\Delta^{-1}(\Delta\sigma_i^{-1}),\qquad
 S\Delta^{-1}=\Delta^{-1}\tau(S),\qquad
 \tau(p)[i]=b-1-p[b-1-i].
\]
After transfers make all pairs left weighted, leading half-twists are
extracted into $z$.  Any half-twist factor must then be initial, since a
left-weighted factor immediately before $\Delta$ must itself be $\Delta$.

\begin{proposition}[Elementary normalization bound]
Normalizing an $m$-letter signed word uses at most $Bm(m-1)/2$ atom
transfers, at most $m$ simultaneous simple factors, and
$O(m^2b^3)$ index-word operations.  Appending one atom to a normal state
with $r$ factors costs $O((r+1)^2b^3)$ conservatively.
\end{proposition}
\begin{proof}
For a temporary factor list indexed from zero, use
$\Phi=\sum_j j\,\ell(s_j)$.  Each transfer decreases it by one before
any deletion.  Deletion and leading half-twist extraction cannot increase
it, and $\tau$ preserves simple length.  The $k$th append raises it by at
most $B(k-1)$ and creates at most one factor.  Summing proves the transfer
bound.  Each pair test or transfer scans $O(b)$ permutation positions;
failed tests advance the scan index, and a successful transfer backs it
up by at most one.  List copying, twisting and deletion fit the stated
quadratic bound.  Starting with $r$ factors instead gives initial
potential $O(B(r+1)^2)$ and the single-append bound.
\end{proof}

The certificate records local transfer positions, generator indices, and
half-twist extraction counts.  A separate verifier imports neither the
normalizer nor the optimizer.  It applies the elementary permutation
operations only after checking their inequalities and indices.  Therefore
it proves equality of the represented braid throughout the trace, even
without trusting normal-form uniqueness or the optimizer's choices.
Version one checks equality of a source interval to the identity or a
signed generator.  Version two checks that source times inverse target
replays to the identity.  It also checks source binding, rotation, disjoint
intervals, strict shortening, target radius, and reconstructed output.
The digest is an input-binding aid; replay on the actual source supplies
the equality proof.  It certifies neither kernel optimality nor a knot
verdict.  A cyclic cut preserves oriented closure by conjugacy, while
uncut substitutions preserve the braid element itself.

\subsection{Sharing the algebra across cuts and targets}

Write $w$ for the input of length $n$, and $P_j$ for the exact prefix
products of the doubled word $ww$.  A source interval $[i,j)$ equals a
target word $v$ precisely when
\[
 P_i^{-1}P_j=v\quad\Longleftrightarrow\quad P_i=P_jv^{-1}.
\]
This ordering matters in a noncommutative group.  The doubled table lets
all $n$ cyclic cuts query the same prefix states.  The radius-$\rho$
dictionary contains every literal signed word of length at most $\rho$,
including the empty word; its size is
\[
 M=\sum_{d=0}^{\rho}(2b-2)^d.
\]
Different words may represent the same braid, which causes no correctness
problem.  A trie of reversed inverse target words shares the appends
needed to evaluate $P_jv^{-1}$ for one endpoint.  Canonical query states
are looked up in the sorted prefix table.

For a fixed cut $s$, let $F[j]$ be the maximum saving on the first $j$
letters.  Immediately before processing $j$, maintain
\[
 H(q)=\max_{0\le i<j:\ P_{s+i}=q}(F[i]-i).
\]
The recurrence is
\[
 F[j]=\max\left(F[j-1],\ \max_{v\in D}
          \{j-|v|+H(P_{s+j}v^{-1})\}\right).
\]
Only afterwards insert $F[j]-j$ under $P_{s+j}$, excluding empty source
intervals.  Backpointers are updated only on a strict improvement.

\begin{proposition}[Optimal disjoint one-pass kernel]
The recurrence returns the shortest word obtainable by disjoint interval
substitutions to the chosen dictionary at that cut.  Minimizing across
all cyclic cuts gives the corresponding cyclic optimum.
\end{proposition}
\begin{proof}
An optimal prefix plan either retains its last letter or ends with a
replacement $[i,j)\mapsto v$.  The latter saves $F[i]+j-i-|v|$ and is
permitted exactly by the prefix-state equality.  Conversely every term
extends a valid optimal prefix plan.  Induction proves the recurrence.
A nonshortening substitution has saving at most $F[i]\le F[j-1]$ and
cannot win strictly, so recorded intervals strictly shorten.  Backtracking
to their starts makes them disjoint.  Any circular disjoint plan admits
a cut at a block boundary or an unused boundary, including a block covering
the entire circle, and is therefore among the tested cuts.
\end{proof}

This is not optimization over arbitrary sequences of substitutions,
conjugations or Markov moves.  It extends the earlier interval dynamic
program by supplying exact unrestricted group queries.  The lower bound
\[
 k\ge\max\left(|\operatorname{exp}(w)|,
                      b-\#\text{cycles of the strand permutation}\right)
\]
allows an exact early exit: each letter changes exponent by one and changes
permutation cycle count by at most one.  Both quantities are unchanged by
conjugacy.  In particular, a knot closure needs at least $b-1$ letters.

\subsection{Height-sensitive costs and the small-core theorem}

Let $R=\max_{j\le2n}\ell_c(P_j)$, where $\ell_c$ counts proper simple
factors.  Inversion preserves this length, and multiplication is
subadditive: move half-twist powers left, concatenate factor lists, and
normalize without increasing their number.  Thus interval prefixes have
height at most $2R$, and target queries at most $R+\rho$; always $R\le2n$.
The following conservative word bound includes a winning proof and replay:
\begin{align*}
 T=O\bigl(&nM(R+\rho+1)^2b^3
       +nM(R+\rho+1)b\log(n+2)\\
       &+n^2M+M\rho\log(M+2)+n(R+\rho+1)b^3\bigr).
\end{align*}
The first term charges the $O(nM)$ exact append queries, the second their
full-state comparisons, and the third all per-cut scalar dynamic programs.
The next term builds the target dictionary and trie.  In the last term,
disjoint winning source intervals total at most $n$ letters and their
strictly shorter targets total at most $n$; the height bound and transfer
potential therefore charge proof generation and replay to
$O(n(R+\rho+1)b^3)$.  Proof size is
$O(n(R+\rho+1)b^2)$ elementary records.  Search workspace, excluding that
output, is
\[
 O\bigl(n(R+1)b+nM+M\rho+\rho(R+\rho+1)b\bigr)
\]
words.  Each index word occupies $O(\log(n+b+\rho+M+2))$ bits, giving
the corresponding bit-storage bound.  Elementary index arithmetic adds at
most a polynomial in this logarithm to the word-operation cost in bit time.  With fixed $\rho$, the computation is polynomial
in $n,b$; leading conservative terms are $O(n^3b^4)$ for radius one and
$O(n^3b^5)$ for radius two.  An arbitrary binary-encoded radius is not a
fixed parameter: its dictionary is exponential in $\rho\log(b+1)$.

A cheaper portfolio stage first expands the whole-word normal form only
if its calculated Artin length is smaller than the input.  It independently
proves this equality.  If the candidate attains the universal lower bound,
it returns immediately; otherwise it compares that candidate with the cyclic
kernel.  The returned word is no longer than the kernel optimum, but need
not itself be an optimal radius-specific interval plan.  Evidence distinguishes
\code{whole-normal-form} from \code{cyclic-kernel}.

\begin{proposition}[Constructive restricted recognition bound]
For fixed radius $\rho$, uncapped certified preprocessing followed by an
uncapped conventional exact recognizer has cost
\[
 \operatorname{poly}_{\rho}(n,b)+\operatorname{poly}(k)2^{O(k)},
\]
where $k$ is the resulting word length.  It is quasi-polynomial on families
whose cyclic radius-$\rho$ optimum is $O(\log^2(n+2))$.
\end{proposition}
\begin{proof}
The optimizer and verifier have the polynomial bounds just given.  Closure
is preserved and remains one component, so $b\le k+1$.  A conventional
Khovanov cube on the resulting $k$ crossings has $2^k$ resolutions with
at most $b+k$ circles each.  Hence its enhanced-state dimension is
$2^{O(k)}$, and exact matrix ranks over $\F$ can be obtained in polynomial
work in that dimension.  Reduced rank one detects the unknot: the rational
unknot-detection theorem and the universal coefficient theorem imply the
same rank-one criterion over $\F$.  The maintained exact fallback supplies
this decision.  Substitution of the small-core bound gives
$2^{O(\log^2(n+2))}$.  Local probe cutoffs are disabled in this theorem.
\end{proof}

There is a proved family for which the decomposition need not be supplied.
Partition a short core into blocks of length at most $\rho$, replace each
by a no-shorter equal braid word, and insert disjoint identity blocks,
forming $y$.  Surround it by $u,v$ with $vu=1$, forming $uyv$.
A cyclic cut exposes $yvu$; deleting the sleeve and identity blocks and
shortening the inflated core blocks is an admissible disjoint plan.
The optimizer discovers a result no longer than the original core without
being told this decomposition.  This generalizes rank-two inflation to
arbitrary supports and multi-letter targets.  It does not establish a small
core for every unknot or a general sub-exponential recognition algorithm.

\subsection{Two limitations of short-word approaches}

First, increasing target radius can be essential even for an unknot.
In $B_4$ write $a=\sigma_1$, $b=\sigma_2$, $c=\sigma_3$ and, for $m\ge2$,
\[
 X_m=a^{-m}c^m a^{m+1}c^{-(m-1)},\qquad W_m=X_m b.
\]
Since $ac=ca$, $X_m=ac$ and the closure of $W_m$ is the unknot.
Nevertheless, the cyclic radius-one optimum is $4m+1$, while radius two
returns three letters.  In $\langle a,c\rangle\cong\mathbb Z^2$, the
prefix path visits successively $(0,0),(-m,0),(-m,m),(1,m),(1,1)$.
It is simple and has no nonconsecutive lattice-neighbor vertices.  Thus
no interval in $X_m$ is an identity or a long signed-letter replacement.
An interval containing the sole $b$ has form $p b q$ with $p,q$ in this
commuting subgroup.  The permutation double cosets force both $p,q$ to
be pure if this equals a signed $b$.  Write $p=a^{2r}c^{2s}$ and
$q=a^{2u}c^{2v}$.  Signed crossings between original strand labels occur
on pairs $12,34,13,24$ with counts $2r,2s,2u,2v$ respectively, and pair
$23$ has count $+1$.  These equality invariants force $p=q=1$ and the
positive sign.  The simple prefix path then permits only the original
single letter, which does not shorten.  This covers intervals at every
cyclic cut.  Radius two replaces $X_m$ by $ac$; the four-cycle permutation
requires at least three output letters.  A whole-word normal form also
exposes the short simple braid, so the practical portfolio avoids the
larger dictionary on this family.  The theorem separates interval rules,
not all braid algorithms.

Second, preserving the braid element itself cannot universally give a
polylogarithmic core, even at three strands.  In $B_3$, the unknot braid
\[
 G_m=a^{m+1}ba^{-m}=a^m(ab)a^{-m}
\]
 has minimum equal-braid Artin length exactly $2m+2$.  To see this, quotient
$B_3$ by its center as $C_2*C_3$, with $x=aba$, $y=ab$.
The image of $G_m$ is $(y^2x)^m y(xy)^m$, of reduced syllable length
$4m+1$.  Each signed Artin letter contributes at most two syllables, so
an equal word has length at least $2m+1$.  Exponent sum two forces even
length, giving $2m+2$, attained by $G_m$.  A shortening subinterval would
contradict this geodesicity.  Cyclic rotation, however, exposes the sleeve
$a^{-m}a^m$ and leaves $ab$ of length two.  Existing free cyclic cleanup
already handles this example.  The obstruction applies to equality-only
explicit words, not to compressed representations or closure-preserving
transformations in general.

\subsection{Production contract and empirical results}

The \code{--garside} option retains the original checked source braid
through PD simplifications.  The proposal compares its verified output
length with the current PD's crossing count, not merely with source length.
An accepted source branch and the earlier PD move trace are recorded
separately; they are not concatenated into a fictitious move sequence.
Restart preserves backend, coefficient, reduction, window, and other user
options, subtracts elapsed time from the global allowance, and disables
this probe to prevent recursive reapplication.

Search, dictionary construction, and independent replay share a cooperative
local deadline and operation allowance.  The host's global check always
runs first.  A local limit or allocation failure skips the probe; global
exhaustion yields \code{UNKNOWN}.  Invalid certificates propagate as errors.
Even the lower-bound shortcut checks resources.  These are cooperative
limits, not hard process-memory or interrupt-latency guarantees; sorting
and individual array operations can be noninterruptible.

The existing modular Alexander obstruction runs before the optional probe.
When it is inconclusive and the probe declines, this result is reused on
the unchanged single-factor diagram.  Individual factors receive their own
checks.  The selected Jones and requested exact Alexander stages follow
the probe: successful compression can avoid their work on the larger PD.
This ordering avoids spending a search allowance on a knot already detected
by the inexpensive modular test, while retaining the opportunity to shorten
before a larger invariant computation.  Source-free PD input bypasses the
probe.  No production default changes.

The delivered suite passes $47$ tests, including $7016$ short-word normal
forms checked against a faithful Artin action and $1512$ independent interval
optimization comparisons.  All $45$ archive hashes and six inspected source
Git blobs were checked.  Three blobs are tied to the named source commit
and path; three were recorded only by exact blob.  The previously unexecuted
adapter smoke now passes.  Production adds normalized-homology comparisons
on $60$ small knot closures, independent replay without the normalizer,
backend/reduction restarts, corrupted certificates, shared resource hooks,
and filter-priority/cache-reuse tests.  All $517$ maintained tests pass.
These are executable checks, not
formal verification of Python semantics or exhaustive knot classification.

\begin{center}
\begin{tabular}{@{}lrrrr@{}}
\toprule
Recognition input & $n$ & Off/radius 1 & Off/radius 2 & A/A \\
\midrule
barrier\_h4 & 51 & 0.979 & 0.971 & 1.003 \\
rectangle\_m16 & 65 & 0.977 & 1.007 & 1.032 \\
geodesic\_m12 & 26 & 0.865 & 0.888 & 0.981 \\
hidden\_sleeve\_m8 & 24 & 0.841 & 0.823 & 1.008 \\
unknot\_h2\_u8 & 43 & 1.018 & 1.037 & 0.989 \\
unknot\_h2\_u16 & 59 & 0.811 & 0.814 & 0.997 \\
figure8\_h2\_u8 & 45 & 1.099 & 1.073 & 0.998 \\
figure8\_h2\_u16 & 61 & 1.102 & 1.107 & 1.014 \\
positive\_h2\_u8 & 49 & 0.775 & 0.763 & 0.961 \\
positive\_h2\_u16 & 65 & 0.834 & 0.827 & 1.018 \\
unknot\_h8\_u8 & 115 & 2.000 & 1.974 & 1.000 \\
unknot\_h8\_u16 & 131 & 2.053 & 2.042 & 0.997 \\
figure8\_h8\_u8 & 117 & 1.094 & 1.099 & 1.005 \\
figure8\_h8\_u16 & 133 & 1.096 & 1.092 & 1.006 \\
positive\_h8\_u8 & 121 & 0.870 & 0.876 & 1.000 \\
positive\_h8\_u16 & 137 & 0.882 & 0.869 & 1.000 \\
unknot\_h16\_u8 & 211 & 1.529 & 1.523 & 1.004 \\
unknot\_h16\_u16 & 227 & 1.340 & 1.327 & 0.988 \\
figure8\_h16\_u8 & 213 & 1.077 & 1.098 & 1.004 \\
figure8\_h16\_u16 & 229 & 1.099 & 1.096 & 0.996 \\
positive\_h16\_u8 & 217 & 0.895 & 0.902 & 0.975 \\
positive\_h16\_u16 & 233 & 0.911 & 0.911 & 0.997 \\
conway & 11 & 0.960 & 0.982 & 0.993 \\
hard\_unknot\_8 & 8 & 0.750 & 0.749 & 0.996 \\
stress\_braid5\_36 & 36 & 1.182 & 1.196 & 1.046 \\
\bottomrule
\end{tabular}
\end{center}
Ratios are medians of paired complete-recognition times; values above one favor the enabled policy.
\begin{center}
\begin{tabular}{@{}lrrr@{}}
\toprule
Isolated compressor & $n$ & Repeated/shared & A/A \\
\midrule
single\_crossing & 1 & 0.891 & 1.004 \\
positive\_4\_braid & 9 & 3.901 & 0.983 \\
ranktwo\_barrier\_h1 & 15 & 7.769 & 1.016 \\
ranktwo\_barrier\_h4 & 51 & 22.510 & 1.000 \\
geodesic\_m12 & 26 & 6.231 & 1.009 \\
hidden\_sleeve\_m8 & 24 & 6.441 & 1.000 \\
seeded\_5\_braid & 22 & 9.832 & 0.979 \\
\bottomrule
\end{tabular}
\end{center}
Both compressor arms search all cuts without lower-bound pruning and generate and replay only the winning certificate.


The complete-recognition experiment keeps the normal filters enabled and
uses fresh validated diagrams, with their construction outside timing.
Probe search, replay, candidate construction, and restarted recognition are
included.  Seven shuffled paired warm rounds compare disabled, identical
control, and radius-one/radius-two enabled policies.  The local allowance is
$0.1$ seconds and $100000$ ticks, with $100000$ target words; all arms have
the same two-second global and $50000$-object ceilings.  An exhausted local
probe is not censored if the remaining exact pipeline completes; a global
\code{UNKNOWN} would receive no speedup ratio.  Raw samples, repeats,
verdicts, stages and source hashes are retained.  Isolated compressor
measurements are reported separately and do not imply the same gain in
recognition.


All $25$ inputs complete with agreeing verdicts in every measured arm.  The
four larger inflated-unknot fixtures have radius-one gains of $1.34$--$2.05$
relative to the disabled pipeline.  The two smaller ones have ratios $1.018$
and $0.811$, so successful shortening need not reduce total work.  Several
inputs are already decided before any probe; enabled-option overhead makes
some very short certificates slower, with a ratio as low as $0.750$.  The
identical-control ratios range from $0.961$ to $1.046$.  These observations
support an optional policy, not a universal or default performance win.
On modular-Alexander-detected examples the enabled path returns before
factorization, which also accounts for modest gains unrelated to Garside
arithmetic.  Source-free PD inputs do not run the probe.

Two retained diagnostic experiments explain the placement choice.  Putting
compression before modular Alexander made some already-detectable knots
roughly $140$ times slower.  Moving all invariant filters before it removed
that regression but lost the useful inflated-unknot gains, because Jones
work on the larger PD had already been spent.  The final placement lies
between modular Alexander and Jones.  These diagnostics retain their own
source hashes and are not pooled with the final timings.

The isolated shared-prefix compressor improves by $3.90$--$22.51$ on the
six multi-crossing fixtures, but loses on the one-letter control.  At the
51-letter barrier, sharing uses $714$ algebra appends instead of $18207$.
That mechanism result is distinct from the full recognition table: the
ordinary pipeline already decides the displayed barrier without invoking
the compressor.  Radius two does not deliver a consistent recognition gain
on this sample, since the inexpensive whole-normal-form proposal often
already attains the lower bound.
